%==============================================================
%  COVER LETTER - master template
%  Change ONLY the two variables below per application:
%    \Position  - the role title (appears twice in the body)
%    \SignDate  - the signature date (DD/MM/YYYY)
%  Everything else stays constant. Build:  tectonic cover_letter_master.tex
%==============================================================
\documentclass[11pt,a4paper]{article}

% Tectonic uses the XeTeX engine. fontspec keeps the PDF text layer ATS-safe:
%  - Ligatures=NoCommon disables fi/fl (they otherwise become single glyphs)
%  - not enabling Ligatures=TeX keeps ' straight (U+0027) and -- as two hyphens
% Requires the Latin Modern OpenType fonts (fonts-lmodern; installed in CI).
\usepackage{fontspec}
% Keeps the PDF text layer ATS-safe under XeTeX/Tectonic:
%  - Mapping=          empties the TeX input mapping, so straight ' and "
%                      stay as U+0027 / U+0022 (no curly quotes)
%  - Ligatures=NoCommon disables fi/fl ligatures (single-glyph parser breakers)
% Requires the Latin Modern OpenType fonts (fonts-lmodern; installed in CI).
\setmainfont{Latin Modern Roman}[Mapping=, Ligatures=NoCommon]
\usepackage[margin=2cm]{geometry}
\usepackage{titlesec}
\usepackage{enumitem}
\usepackage{xcolor}
\usepackage{hyperref}
\usepackage{parskip}

% ATS-safe list marker: a plain hyphen, not a bullet glyph.
\renewcommand{\labelitemi}{-}

\definecolor{accent}{HTML}{2C3E50}
\hypersetup{colorlinks=true, urlcolor=accent, linkcolor=accent}

\titleformat{\section}{\large\bfseries\color{accent}}{}{0em}{}
\titlespacing{\section}{0pt}{10pt}{4pt}
\setlist[itemize]{leftmargin=1.4em, itemsep=3pt, topsep=3pt}
\pagestyle{empty}

% ---- PER-APPLICATION VARIABLES (edit these two only) --------
\newcommand{\Position}{Principal C\&SP Methodology Lead}
\newcommand{\SignDate}{05/10/2026}
% ------------------------------------------------------------

\begin{document}

{\Large\bfseries\color{accent} Valentin Legras}\\[10pt]

\section{About myself}
As a Senior Data Scientist and Product Owner with 7+ years in regulated clinical
data science, my goal is to secure the \Position{} position. I work across
clinical and statistical programming workflows using SAS and R, with current
hands-on work split about 40\% SAS and 60\% R. My experience covers CDISC
SDTM/ADaM datasets, Pinnacle 21 checks, GxP software practices, regulatory
submission support, data anonymisation, and clinical data privacy. I have also
led the 'aNCA' open-source R package from concept to enterprise-wide adoption,
with 120+ users and more than 2M CHF in annual savings.

\section{My mindset}
Strategy and automation entered my life long before I knew the words for them. I
was around 10 years old when I started competing in basketball. To perform well
against top teams across Europe, individual skill was never enough. It required
strategy: analyzing opponents, optimizing plays, and finding the most efficient
way to win.

Around the same time, I dived into strategic video games. I was driven to compete
in environments where others could pay to get bonuses. To succeed, I had to find
the most optimized path, meticulously planning and automating processes to
maximize outcomes with limited resources. This wasn't just playing; it was a
challenge to reverse-engineer systems, think out of the box, think about how
other brains work and find the most efficient solution.

This mindset, which revolves around strategy, optimization, and automation, is a
core part of my identity. I have a strong drive to find the most efficient and
strategic solution to any problem. In the future, I view AI as a transformative
tool that amplifies this drive, enabling me to automate complex processes and
uncover optimization opportunities at a scale that wasn't possible before (via
multi-agent coding). I don't only build robust and optimized systems in my
professional life but also for my personal interests. As an example, I developed
back-testing reproducible statistical algorithms to optimize strategies on sports
betting websites and investment on S\&P500 stocks (see GitHub).

Having successfully led, coached, and mentored the 'aNCA' package team of 7
developers and onboarded new clinical and pre-clinical colleagues, I've found my
greatest impact comes from combining this strategic approach with team
leadership. My vision is to apply this lifelong passion for optimization at a
larger scale, leading teams to build the most efficient, automated, and impactful
solutions possible.

I would like to explain why my current profile fits the \Position{} role. My
background is rooted in PK/PD and pharmacometrics, but my day-to-day work is not
limited to PK data. I handle clinical data from many domains, including
time-to-event, adverse events, laboratory, efficacy, dose, exposure, and safety
data. This gives me a practical view of how reusable tools, standards, and
automation need to work across functions rather than only inside one specialty.

\begin{itemize}
  \item \textbf{Clinical programming methodology:} I have converted legacy SAS
        processes into reproducible R pipelines, automated ADPC, ADNCA, and ADPP
        dataset generation, and built tools used by clinical and pre-clinical
        scientists at scale.
  \item \textbf{SAS and R expertise:} I am comfortable in both SAS
        (BASE/MACRO/SQL) and R, and I understand when a regulated clinical
        programming workflow needs one stack, the other, or a controlled bridge
        between them.
  \item \textbf{Regulated delivery:} My work includes GxP software standards,
        testing, documentation, code review, QC, Pinnacle 21, health authority
        responses, electronic submission support, and careful handling of
        confidential or anonymised clinical data.
  \item \textbf{Leadership and adoption:} As Product Owner of 'aNCA', I led and
        mentored 7 developers, onboarded new users, managed a 15+ stakeholder
        network, and turned a manual workflow into a 90\% automated process.
  \item \textbf{Innovation with control:} I scale AI-assisted coding workflows
        for issue creation, automated problem solving, and code review, while
        keeping engineering standards, traceability, and review discipline in
        place.
\end{itemize}

I am ready to apply this combination of SAS/R expertise, clinical data standards,
tool development, and practical leadership to build methodology solutions that
increase automation, quality, and consistency across clinical and statistical
programming teams.

\vspace{10pt}
Sincerely\\[6pt]
Valentin Legras\\
\SignDate{}

\end{document}
